\chapter{Cycles, inversions, and the major index}
\label{ch:inversions}

This chapter formalizes Stanley EC1 \S1.3 (cycles and inversions), the
section that introduces the two foundational permutation statistics
beyond descents: the \emph{cycle count} $c(w)$ and the two Mahonian
statistics \emph{inv} and \emph{maj}.  The headline result of the
section --- the classical \emph{inv-maj equidistribution} of MacMahon
--- is part of an in-progress extension; the present chapter covers the
definitions and the easy identities, with explicit pointers to what
remains.

\medskip
This chapter uses the same indexing convention as
Chapter~\ref{ch:descents}: Stanley's $w \in \mathfrak{S}_n$ becomes our
$\sigma \in \Snp{n+1}$, and Stanley's descent positions $1, \dots, n-1$
become our $\{0, \dots, n-1\}$ via the index shift.

% ============================================================
\section{Cycle representation of permutations}
% ============================================================

\begin{definition}[Cycle count]
  \label{def:cycle_count}
  \rocq{mathcomp_eulerian.cycles.cycle_count}
  \rocqok
  For a permutation $\sigma \in \Snp{n}$, the \emph{cycle count}
  $c(\sigma)$ is the number of cycles in the disjoint-cycle
  decomposition of $\sigma$.  In MathComp this is the cardinality of
  $\mathrm{porbits}(\sigma)$, the set of orbits of $\sigma$ acting on
  $\{0, \dots, n-1\}$ by iteration.
\end{definition}

\begin{lemma}[$c(\sigma) \leq n$]
  \label{lem:cycle_count_le}
  \uses{def:cycle_count}
  \rocq{mathcomp_eulerian.cycles.cycle_count_le}
  \rocqok
  Every cycle is non-empty, so the cycle count is bounded by $n$.
\end{lemma}

\begin{lemma}[Identity has $n$ cycles]
  \label{lem:cycle_count_id}
  \uses{def:cycle_count}
  \rocq{mathcomp_eulerian.cycles.cycle_count_id}
  \rocqok
  Each fixed point is a 1-cycle, so $c(\mathrm{id}) = n$.
\end{lemma}

% ============================================================
\section{Stirling numbers of the first kind}
% ============================================================

\begin{definition}[Signless Stirling cycle numbers]
  \label{def:stirling_c}
  \uses{def:cycle_count}
  \rocq{mathcomp_eulerian.cycles.stirling_c}
  \rocqok
  $c(n, k) = \#\{\sigma \in \Snp{n} : c(\sigma) = k\}$ counts
  permutations of $\{0, \dots, n-1\}$ with exactly $k$ cycles in their
  disjoint-cycle decomposition.  Stanley writes $c(n, k)$ throughout
  \S1.3.2.
\end{definition}

\begin{lemma}[Row sum]
  \label{lem:stirling_c_row_sum_fact}
  \uses{def:stirling_c}
  \rocq{mathcomp_eulerian.cycles.stirling_c_row_sum_fact}
  \rocqok
  $\sum_{k=0}^{n} c(n, k) = n!$.  Each permutation is counted by
  exactly one $k$, so the row sum recovers the total.
\end{lemma}

\begin{theorem}[Stirling recurrence]
  \label{thm:stirling_c_rec}
  \uses{def:stirling_c}
  \rocq{mathcomp_eulerian.stirling_fiber.stirling_c_rec}
  \rocqok
  For all $n, k \in \NN$,
  \[
    c(n+1, k+1) \;=\; n \cdot c(n, k+1) + c(n, k).
  \]
  Combinatorially: $n+1$ either becomes a fixed point of a
  $k$-cycle permutation (giving $c(n, k)$), or is inserted into one
  of $n$ positions of a $(k+1)$-cycle permutation (giving
  $n \cdot c(n, k+1)$).

  \emph{Formalization route.}  The proof factors through a
  construction $\mathrm{insert\_after}\,j_0\,s : \Snp{n+1}$ that
  splices $\mathrm{ord\_max}$ into the $s$-cycle through $j_0$.
  The key algebraic identity
  \[
    \mathrm{insert\_after}\,j_0\,s
    \;=\;
    \mathrm{tperm}\,(\mathrm{ord\_max})\,(\mathrm{lift}\,\mathrm{ord\_max}\,j_0)
    \;\cdot\;
    \mathrm{lift\_perm}\,\mathrm{ord\_max}\,\mathrm{ord\_max}\,s
  \]
  reduces cycle-count preservation to MathComp's
  \texttt{porbits\_mul\_tperm} composed with the fixed-point
  cycle-count lemma in
  \rocq{mathcomp_eulerian.cycles_rec.cycle_count_lift_perm_id}.
\end{theorem}

% ============================================================
\section{Inversions}
% ============================================================

\begin{definition}[Inversion at a pair]
  \label{def:is_inv}
  \rocq{mathcomp_eulerian.inversions.is_inv}
  \rocqok
  Position pair $(i, j)$ is an \emph{inversion} of $\sigma$ when
  $i < j$ but $\sigma(j) < \sigma(i)$ — i.e., the pair appears in the
  ``wrong'' order.
\end{definition}

\begin{definition}[Inversion set and count]
  \label{def:inv}
  \uses{def:is_inv}
  \rocq{mathcomp_eulerian.inversions.inv_set}
  \rocq{mathcomp_eulerian.inversions.inv}
  \rocqok
  $\mathrm{inv\_set}(\sigma)$ is the set of inversion pairs;
  $\mathrm{inv}(\sigma) = |\mathrm{inv\_set}(\sigma)|$ is their count.
  Stanley writes $\mathrm{inv}(w)$.
\end{definition}

\begin{lemma}[Identity has no inversions]
  \label{lem:inv_id}
  \uses{def:inv}
  \rocq{mathcomp_eulerian.inversions.inv_id}
  \rocqok
  $\mathrm{inv}(\mathrm{id}) = 0$.
\end{lemma}

% ============================================================
\section{The major index}
% ============================================================

\begin{definition}[Major index]
  \label{def:maj}
  \uses{def:descent_set}
  \rocq{mathcomp_eulerian.inversions.maj}
  \rocqok
  The \emph{major index} is the sum of (1-indexed) descent positions:
  \[
    \mathrm{maj}(\sigma) \;=\; \sum_{i \in \DescSet{\sigma}} (i + 1)
    \;=\; \sum_{i \in D(w)} i \quad (\text{Stanley's notation, p.~22}).
  \]
\end{definition}

\begin{lemma}[Identity has zero major index]
  \label{lem:maj_id}
  \uses{def:maj}
  \rocq{mathcomp_eulerian.inversions.maj_id}
  \rocqok
  $\mathrm{maj}(\mathrm{id}) = 0$.
\end{lemma}

% ============================================================
\section{The classical equidistribution — work in progress}
% ============================================================

\begin{remark}[Inv-maj equidistribution — not yet formalized]
  \label{rem:inv_maj_equidistr}
  The headline of Stanley \S1.3.3, MacMahon's equidistribution result,
  is:
  \[
    \forall k,\quad
    \#\{\sigma \in \Snp{n+1} : \mathrm{inv}(\sigma) = k\}
    \;=\;
    \#\{\sigma \in \Snp{n+1} : \mathrm{maj}(\sigma) = k\}.
  \]
  Foata's first fundamental bijection $\varphi : \Snp{n+1} \to
  \Snp{n+1}$ satisfies $\mathrm{inv}(w) = \mathrm{maj}(\varphi(w))$
  and witnesses the equidistribution constructively.  Both sides are
  generated by the q-factorial
  $[n+1]_q! = \prod_{k=1}^{n+1} (1 + q + \dots + q^{k-1})$.

  This result will land in a new file \texttt{foata.v} once Phase 3 of
  the §1.3 extension plan completes.  See
  \texttt{docs/plans/INVERSIONS\_PLAN.md} for the planned scope.
\end{remark}
